\section{Numerical experiments}\label{sec:experiments}

We address three questions: how MIST compares with CPU interior-point solvers at
the same or higher accuracy (\cref{sec:pushbox,sec:pendulum}); how much the
hierarchy gains over a general sparse factorization on the same GPU
(\cref{sec:cudss}); and whether the time is spent where the analysis predicts
(\cref{sec:phases}).

\subsection{Setup}\label{sec:setup}

\paragraph{Problems.}
PushT and PushBox are contact-rich planar pushing tasks from SPOT~\cite{kang2025spot};
the third problem is a three-dimensional pendulum maneuver. All are relaxed at
order $r=2$ and exported with stage labels; the largest PushT instance ($N=100$)
has $259{,}499$ equations and $63{,}885$ Gram blocks. \todo{table of instance
sizes ($m$, number and orders of Gram blocks) for all problems.}

\paragraph{Solvers and hardware.}
We compare MIST with MOSEK~11.0.10~\cite{mosek} and BRISK~1.6.5~\cite{brisk} on
the same workstation (one H100 for MIST, CPU for the baselines), and with
full-system cuDSS~0.7.1.6~\cite{cudss} inside the same IPM on one A100.
\todo{CPU model and thread counts.}

\paragraph{Accuracy and timing.}
Every returned solution is checked independently on the original SDP using eight
relative metrics: primal and dual equality residuals, duality gap,
complementarity, the distances of the primal and dual matrices to the PSD cone,
and their symmetry residuals. A solution passes the \emph{standard} check if all
metrics are below $10^{-4}$, and the \emph{strict} check if all are below
$10^{-6}$ with unaveraged complementarity. Times run from reading the problem to
writing the solution; the independent check is not timed.

\subsection{Comparison with CPU solvers}\label{sec:pushbox}

\Cref{tab:pushbox,fig:pushbox} report PushBox for horizons from 2 to 80. MIST
selects its CPU backend for $N\le5$ and the GPU backend from $N=10$ on.

\begin{table}[htbp]
\centering
\caption{PushBox: time (s, median over repeated runs) and number of runs whose
solution passes the strict $10^{-6}$ check on the original SDP.}
\label{tab:pushbox}
\small
\begin{tabular}{@{}r rrr ccc@{}}
\toprule
& \multicolumn{3}{c}{Time (s)} & \multicolumn{3}{c}{Strict check passed}\\
\cmidrule(lr){2-4}\cmidrule(l){5-7}
$N$ & MIST & BRISK & MOSEK & MIST & BRISK & MOSEK\\
\midrule
2  & 0.70 & 0.62 & 0.73 & 1/1 & 2/2 & 2/2\\
5  & 2.48 & 3.61 & 2.38 & 3/3 & 0/2 & 0/2\\
10 & 4.13 & 4.84 & 3.63 & 3/3 & 0/2 & 0/2\\
20 & \textbf{5.46} & 13.21 & 8.42 & 3/3 & 0/1 & 0/1\\
50 & \textbf{7.32} & 13.23 & 20.89 & 3/3 & 0/1 & 0/1\\
80 & \textbf{9.20} & 29.34 & 28.61 & 3/3 & 0/1 & 0/1\\
\bottomrule
\end{tabular}
\end{table}

\begin{figure}[htbp]
\centering
\begin{tikzpicture}
\begin{loglogaxis}[
  width=.68\linewidth,height=5.8cm,
  xlabel={Horizon $N$},ylabel={Time (s)},
  xmin=1.7,xmax=95,ymin=0.4,ymax=45,
  xtick={2,5,10,20,50,80},xticklabels={2,5,10,20,50,80},
  ytick={0.5,1,2,5,10,20,40},yticklabels={0.5,1,2,5,10,20,40},
  grid=both,grid style={gray!18},
  legend style={at={(0.03,0.97)},anchor=north west,draw=none,font=\small,fill=white,fill opacity=0.85,text opacity=1},
  legend cell align=left]
\addplot[paperblue,very thick,mark=*,mark size=2.4pt]
  table[x=N,y=gipm_sec,col sep=comma]{data/pushbox_strict.csv};
\addlegendentry{MIST}
\addplot[paperorange,thick,mark=square*,mark size=2.2pt,mark options={fill=white}]
  table[x=N,y=mosek_sec,col sep=comma]{data/pushbox_strict.csv};
\addlegendentry{MOSEK 11}
\addplot[paperteal,thick,densely dashed,mark=triangle*,mark size=2.8pt,mark options={fill=white,solid}]
  table[x=N,y=brisk_sec,col sep=comma]{data/pushbox_strict.csv};
\addlegendentry{BRISK 1.6.5}
\addplot[paperorange,only marks,mark=square*,mark size=2.2pt,forget plot] coordinates {(2,0.73)};
\addplot[paperteal,only marks,mark=triangle*,mark size=2.8pt,forget plot] coordinates {(2,0.62)};
\end{loglogaxis}
\end{tikzpicture}
\caption{PushBox time versus horizon. Filled markers: all runs pass the strict
check; hollow markers: none does.}
\label{fig:pushbox}
\end{figure}

For $N\le10$ the three solvers are within a factor of about 1.5, since these
instances are too small to occupy a GPU. From $N=20$ on, MIST is $1.5$--$3.1\times$
faster than MOSEK and $1.8$--$3.2\times$ faster than BRISK, and its time grows by
only $1.7\times$ from $N=20$ to $N=80$, against $3.4\times$ for MOSEK. MIST's
solutions pass the strict check in every run, whereas for $N\ge5$ no MOSEK or
BRISK solution does: MOSEK stops with its own optimality status but leaves primal
residuals and complementarity above $10^{-6}$ on the original SDP, and BRISK
reports numerical difficulties from $N=10$ on.

\subsection{A long-horizon instance: 3D pendulum}\label{sec:pendulum}

\Cref{tab:pendulum} reports the three-dimensional pendulum with $N=20$, solved by
MOSEK with eight CPU threads and by MIST on one H100.

\begin{table}[htbp]
\centering
\caption{3D pendulum, $N=20$.}
\label{tab:pendulum}
\small
\begin{tabular}{@{}lrrc@{}}
\toprule
Solver & Iterations & Solver time (s) & Converged (strict check)\\
\midrule
MOSEK (8 CPU threads) & 21 & 1{,}173.0 & yes\\
MIST (H100) & 35 & 95.1 & yes\\
\bottomrule
\end{tabular}
\end{table}

Both solvers converge. MIST takes more iterations, but each costs $2.7$\,s instead of $55.9$\,s, so it is $12.3\times$ faster overall.

\subsection{Hierarchical elimination versus general sparse factorization}\label{sec:cudss}

To isolate the hierarchy, we run the same IPM on the same A100 with two linear
backends: full-system cuDSS on the assembled $H_\delta$, and the hierarchical
backend of \cref{sec:algorithm}. Each backend ran one warm-up and three measured
runs per horizon on PushT; all runs pass the standard check, and the hierarchical
backend never fell back to cuDSS.

\begin{table}[htbp]
\centering
\caption{PushT on one A100: solver time, median [min, max] of three runs, with the
same IPM and two linear backends.}
\label{tab:cudss}
\small
\begin{tabular}{@{}rllcl@{}}
\toprule
$N$ & Full cuDSS (s) & Hierarchical (s) & Speedup & Reduced solve\\
\midrule
10  & 1.455 [1.443, 1.460] & 1.446 [1.444, 1.457] & $1.01\times$ & dense\\
30  & 3.124 [3.099, 3.125] & 2.629 [2.548, 2.633] & $1.19\times$ & dense\\
100 & 9.276 [9.134, 9.329] & 7.624 [7.612, 7.638] & $1.22\times$ & cyclic reduction\\
\bottomrule
\end{tabular}
\end{table}

The two backends tie at $N=10$, and the hierarchy is $1.19\times$ and
$1.22\times$ faster at $N=30$ and $N=100$, with similar iteration counts. These
ratios include assembly and all other IPM work, which both backends share. Since
cuDSS could in principle find an equally good ordering (\cref{prop:ordering}),
the gain comes from execution: a small reduced system, batched factorization of
identical stages, and no global sparse analysis.

\subsection{Where the time goes}\label{sec:phases}

\Cref{tab:phases} breaks down one PushT run with $N=100$ on an A100 with the
current version of MIST.

\begin{table}[htbp]
\centering
\caption{Phase times for PushT, $N=100$, one A100. Indented rows are parts of the
row above.}
\label{tab:phases}
\small
\begin{tabular}{@{}lrr@{}}
\toprule
Phase & Time (ms) & Share of Newton loop\\
\midrule
Setup (assembly structures, interfaces, symbolic plans) & 1{,}373 & --\\
Newton loop, 24 iterations & 1{,}748 & 100\%\\
\quad Local partial factorization & 618 & 35\%\\
\quad Forward condensation & 303 & 17\%\\
\quad Interior recovery & 137 & 8\%\\
\quad Reduced factorization and solves & 21.5 & 1.2\%\\
\bottomrule
\end{tabular}
\end{table}

An iteration takes about 73\,ms. The reduced system costs about $1\%$ of the
Newton loop, while local factorization, condensation and recovery dominate. This
is what \cref{thm:main} predicts once pruning and coarsening have made the
interfaces small: the cost is governed by the local front terms $\Phi_2$ and
$\Phi_1$, and that is exactly the part MIST executes as batched GPU kernels.
